\section{Moteur 3 --- Mean Reversion Journalière par RSI (10\,\%)}
\label{sec:rsi_daily}

\subsection{Intuition économique}

\begin{thesisbox}
\textbf{Un mouvement directionnel exacerbé sur une journée est un \emph{overshoot} statistiquement rare.} Lorsque l'indicateur RSI journalier d'une paire FX majeure dépasse 75 (sur-achat) ou descend sous 25 (sur-vente), il reflète une accumulation disproportionnée de flux dans une seule direction~--- généralement déclenchée par une annonce macroéconomique ou un \emph{squeeze} technique. L'expérience empirique montre que ces extrêmes se résolvent typiquement en 2 à 10 jours par un retour du prix vers un RSI neutre autour de~50.
\end{thesisbox}

Ce moteur exploite \emph{exactement le régime opposé} à celui du Moteur~2~: là où le suivi de tendance capture les mouvements lents déployés sur plusieurs semaines, le RSI journalier capture les excès de court terme. D'où la corrélation négative observée ($-0.251$ sur la période complète) entre ces deux composantes, qui est le levier central de la diversification du portefeuille.

\subsection{Mécanique du signal}

\begin{enumerate}
\item \textbf{Calcul du RSI(14)} sur le close journalier de chaque paire.
\item \textbf{Règles d'entrée}~:
  \begin{itemize}
  \item Long si $\mathrm{RSI} < 25$ (sur-vente)~;
  \item Short si $\mathrm{RSI} > 75$ (sur-achat).
  \end{itemize}
\item \textbf{Règle de sortie}~: retour à un niveau neutre $\mathrm{RSI} \approx 50$ ou inversion du signal.
\item \textbf{Portefeuille du moteur}~: moyenne \emph{equal-weight} des quatre paires EUR-USD, GBP-USD, USD-JPY, USD-CAD.
\end{enumerate}

\textit{Remarque importante}~: contrairement au Moteur~2, USD-CAD est \textbf{conservé} ici. La microstructure \emph{range-bound} qui handicape le momentum sur cette paire est précisément ce qui avantage un signal de retour à la moyenne. Un même défaut peut être une qualité dans un autre contexte stratégique.

\subsection{Le paradoxe du Sharpe standalone}

\begin{warningbox}
\textbf{Le Moteur~3 a un Sharpe standalone de seulement 0.16 sur la période complète.} C'est le plus faible des trois. Pourtant, son retrait ferait perdre environ 30\,\% du Sharpe du portefeuille combiné~--- la définition exacte d'un diversificateur anti-corrélé de valeur.
\end{warningbox}

La décomposition année par année révèle le mécanisme~:

\begin{center}
\renewcommand{\arraystretch}{1.1}
\begin{tabular}{lrrr p{4cm}}
\toprule
\textbf{Année} & \textbf{Moteur 3} & \textbf{Moteur 1} & \textbf{Moteur 2} & \textbf{Commentaire} \\
\midrule
2019 & \textcolor{rsiGreen}{$+0.95\,\star$} & $-0.64$ & $-0.41$ & Le Moteur~3 sauve l'année \\
2020 & $-0.03$ & $+1.19$ & $-0.24$ & --- \\
2021 & $+0.61$ & $0.00$ & $+0.62$ & --- \\
2022 & $-0.51$ & \textcolor{rsiGreen}{$+2.23$} & $+1.18$ & --- \\
2023 & \textcolor{rsiGreen}{$+0.92\,\star$} & $-0.28$ & $+0.55$ & Le Moteur~3 sauve l'année \\
2024 & $-0.24$ & $+2.08$ & $+0.61$ & --- \\
2025 & $+0.38$ & $+0.72$ & $+1.57$ & --- \\
2026 YTD & \textcolor{rsiGreen}{$+3.54\,\star$} & $+1.72$ & $-1.59$ & Le Moteur~3 sauve le YTD \\
\bottomrule
\end{tabular}
\end{center}

\textbf{Le Moteur~3 est positif exactement dans les années où les deux autres perdent.} C'est un diversificateur anti-corrélé au sens strict. La matrice de corrélation sur la période complète confirme~:

\[
\begin{array}{l|rrr}
 & \text{Moteur 1} & \text{Moteur 2} & \text{Moteur 3} \\
\hline
\text{Moteur 1} & 1.000  & +0.056 & -0.027 \\
\text{Moteur 2} & +0.056 & 1.000  & \mathbf{-0.251} \\
\text{Moteur 3} & -0.027 & \mathbf{-0.251} & 1.000 \\
\end{array}
\]

\subsection{Performance historique}

\begin{figure}[H]
\centering
\includegraphics[width=\linewidth]{sleeve_rsi_daily_4p_equity.png}
\caption{Équité et drawdown du Moteur~3 (quatre paires) pris isolément (non pondéré, non leveragé~--- avant application des poids $80/10/10$ et du ciblage de volatilité global). La courbe est relativement plate, mais les pics positifs en 2019, 2023 et début 2026 correspondent exactement aux années faibles des autres moteurs~--- c'est le profil recherché pour un diversificateur.}
\end{figure}

\subsection{Épisodes représentatifs}

\begin{table}[H]
\centering
\caption{Épisodes représentatifs du sleeve RSI Daily 3-pair. Les épisodes courts typiques (2-10 jours) reflètent les retours rapides après sur-achat/sur-vente capturés par RSI(14).}
\label{tab:trades_rsi_daily}
\begin{tabular}{llcrr}
\toprule
Entrée & Sortie & Sens & Durée & Rendement \\
\midrule
2020-03-24 & 2020-03-27 & Long & 4 j & 48.49\,\% \\
2022-09-28 & 2022-09-29 & Long & 2 j & 29.36\,\% \\
2019-05-22 & 2019-05-23 & Short & 2 j & -0.18\,\% \\
2022-09-20 & 2022-09-23 & Short & 4 j & -19.72\,\% \\
2020-03-16 & 2020-03-18 & Short & 3 j & -46.73\,\% \\
\bottomrule
\end{tabular}
\end{table}


\begin{insightbox}
\textbf{La leçon de construction de portefeuille~:} un moteur à Sharpe individuel faible peut être le meilleur diversificateur si sa corrélation avec le reste est négative sur les années difficiles. Ce qui compte pour le portefeuille global n'est pas le Sharpe individuel mais la \textbf{contribution marginale au Sharpe combiné}, qui dépend directement de la corrélation. Un signal à 0.16 anti-corrélé bat couramment un signal à 0.60 fortement corrélé.
\end{insightbox}
